\documentclass{llncs}
%
\usepackage[T1]{fontenc}
\usepackage[utf8]{inputenc}
\usepackage{graphicx}
\usepackage{amsmath}
\usepackage{amssymb}
\setlength{\fboxsep}{6pt}
\newsavebox{\editorcommentbox}
\newenvironment{editorcomment}
    {\begin{lrbox}{\editorcommentbox}\begin{minipage}{0.95\linewidth}\small}
    {\end{minipage}\end{lrbox}\par\medskip\noindent\fbox{\usebox{\editorcommentbox}}\par\medskip}
%
\begin{document}
%
\title{Diagnosing Causal Reasoning Failures in Large Language Models: Implications for LLM-Assisted Causal Discovery}
%
\titlerunning{Diagnosing Causal Reasoning Failures in LLMs}
%
\author{Laura Itzel Rodríguez Dimayuga\inst{1}\orcidID{0009-0007-0122-5294} \and Enrique Francisco Soto Astorga\inst{2}\orcidID{0000-0002-1270-8706}}
%
\authorrunning{L. I. Rodríguez Dimayuga and E. F. Soto Astorga}
%
\institute{Facultad de Ciencias, Universidad Nacional Autónoma de México, Ciudad de México, México\\
\email{laudima@ciencias.unam.mx} \and
Facultad de Ciencias, Universidad Nacional Autónoma de México, Ciudad de México, México\\
\email{astorga@ciencias.unam.mx}}
%
\maketitle
%
\begin{abstract}
Causal discovery methods increasingly rely on background knowledge to compensate for the fact that 
observational data alone cannot, in general, identify a unique causal structure. Large language models (LLMs) 
are progressively used as a convenient source of such background knowledge, proposing causal edges, priors, or 
full graphs from variable names and textual descriptions. This growing practice presupposes that LLMs possess 
causal competence rather than sophisticated pattern completion, an assumption that remains empirically and theoretically underexamined. We present a preliminary, proof-of-concept framework for diagnosing whether LLM causal judgments reflect genuine causal reasoning or associative mimicry. The framework integrates Judea Pearl's Ladder of Causation, which distinguishes associational (L1), interventional (L2), and counterfactual (L3) reasoning, with Robert Rosen's modeling relation, which requires that a legitimate model preserve the structural linkages of the natural system it represents. We derive from this integration a falsifiable prediction: if LLMs operate predominantly at L1, their failures on causal benchmarks should concentrate structurally at L2 and L3 rather than being uniformly distributed. We describe an evaluation protocol, built on the CLadder benchmark, that queries LLMs across causal levels and visualizes failure clustering as a directed acyclic graph (DAG) of failure co-occurrence. We report on the current, pre-empirical stage of this work and discuss its relevance for practitioners who use LLMs to elicit causal priors in hybrid causal discovery pipelines.

\keywords{Causal discovery \and Large language models \and Causal reasoning \and Pearl's causal hierarchy \and Benchmark evaluation \and CLadder \and Rosen's modeling relation}
\end{abstract}
%
\section{Introduction and Motivation}


Observational data alone does not, in general, identify a unique causal 
structure. Distinct DAGs can be Markov equivalent and thus indistinguishable 
from the observed distribution \cite{spirtes2000causation,glymour2019review,verma1990equivalence}. Discovery algorithms address this non-identifiability 
with interventions, temporal information, or background knowledge that orients
edges and prunes the equivalence class \cite{spirtes2016causal}. Such 
background knowledge has traditionally come from human domain experts; 
recently, LLMs have been proposed as a scalable substitute, queried for 
pairwise orientation, edge priors, or full causal graphs from variable names 
and textual metadata, with their outputs fed directly into discovery pipelines 
\cite{long2023can,kiciman2023causal,verma2025causal,butkus2025causal}. This 
practice silently presupposes that LLMs possess causal competence comparable to 
a domain expert. Though some report encouraging 
benchmark performance \cite{kiciman2023causal}, while others argue that
apparent competence is explainable by surface-level statistical association
rather than genuine causal inference \cite{Zecevic2023CausalPL,willig2022can,li2023counterfactual,chen2024clear,hobbhahn2022investigating}. Rather than 
resolve this dispute empirically, we contribute a theoretically motivated 
diagnostic framework and evaluation protocol aimed at determining \emph{where}, structurally, LLM causal reasoning fails, making the epistemic risk of LLM-assisted causal discovery explicit and testable.

\begin{editorcomment}
Comment from editor: In the Introduction, when discussing ways to address non-identifiability, the authors forgot to mention non-Gaussian distributions exploited by approaches such as FASK or ICA-based ones like LiNGAM.
\end{editorcomment}

\section{Theoretical Framework}

Pearl organizes causal reasoning into three layers of increasing informational demand \cite{pearl2009causality,pearl2018bookOfWhy}: \emph{association} (L1), regularities of $P(y \mid x)$; \emph{intervention} (L2), the effect of actively setting a variable, $P(y \mid do(x))$, which requires knowledge beyond any observational distribution; and \emph{counterfactual} (L3), reasoning about a specific unit under a hypothetical alternative, $P(y_x \mid x', y')$, which presupposes a fully specified structural causal model. No amount of L1 information determines L2 or L3 quantities without additional structural assumptions \cite{pearl2018theoretical}: a system could \emph{appear} causally competent while operating strictly at L1, reproducing correct associational answers while lacking the structural information L2/L3 queries require.

Independently, Rosen's theory of anticipatory systems locates causality not in a primitive property of events but in the \emph{linkages} between the observable qualities of a natural system, expressed as its laws of state change \cite{rosen2012anticipatory}. A model is causally legitimate only insofar as a \emph{modeling relation} holds: a structure-preserving correspondence between inferential entailment in the formal system and causal entailment in the natural system it represents. A model that reproduces surface observables without preserving this correspondence is not a causal model of the system, regardless of predictive accuracy.

Together, the frameworks conform a theoretical foundation for diagnosing the limitations of LLM causal reasoning. Pearl's hierarchy specifies \emph{what} information a causal system must encode beyond association; Rosen's modeling relation specifies \emph{why} a system that merely mimics correct L1 answers cannot be expected to satisfy L2/L3 queries reliably, since doing so is not underwritten by any structural correspondence with the system it describes. This yields the testable prediction organizing our protocol: \emph{if an LLM's causal competence is L1-bound, its errors on a stratified causal benchmark should be structurally concentrated at L2 and L3, not uniformly distributed across levels and subtypes.} Uniform errors would instead suggest generic task difficulty, weakening the associative-mimicry account.

\begin{editorcomment}
Comment from editor: I am not familiar with Rosen’s framework, but the way it is described in the second paragraph of part 2 seems to be related to the concepts of Markov and Faithfulness assumptions for causal graphical models. Do the authors agree with this? If so, maybe those assumptions also can provide theoretical foundations to the approach suggested here.
\end{editorcomment}

\begin{editorcomment}
Comment from editor: I do not see, apart from a seemingly vague aspect of Rosen’s framework relating models and nature, what are the authors leveraging from this framework for their own approach. Can your framework work without Rosen’s? I am confused about the role of Rosen’s framework in this project.  
\end{editorcomment}

How to connect the causal relations of Rosen as a theorical statement to frame why LLMs are not able to reason causally.


\section{Evaluation Protocol and Proof of Concept}

We operationalize this criterion using CLadder \cite{jin2023cladder}, roughly 10,000 items generated from causal graphs paired with associational, interventional, and counterfactual queries, each with a ground-truth answer from a symbolic oracle. Unlike naturalistic-text benchmarks, CLadder's ground truth is tied directly to an explicit graph and query type, letting failures be attributed to a specific hierarchy level and graph topology rather than to question ambiguity.

The planned pipeline: (i) sample a stratified subset of CLadder across levels (L1/L2/L3) and subtypes; (ii) submit each item to one or more LLMs via API under a fixed prompting protocol; (iii) score responses against the oracle and log level, subtype, and graph topology; (iv) aggregate failures into a directed graph, built with \texttt{networkx}, whose nodes represent (level, subtype, topology) strata, used to visualize whether errors cluster at L2/L3 or spread diffusely. This failure DAG exposes the \emph{shape} of the model's error surface rather than recovering any object-level causal structure.



\begin{editorcomment}
Comment from editor:  - In Section 3, what do the authors refer to by “subtypes” and “graph topology”.
\end{editorcomment}
\begin{editorcomment}
Comment from editor: model’s error surface
\end{editorcomment}

This are parameters that are used to generate the queries, and they are used to analyze the results. The subtypes are the different types of queries that can be generated from a causal graph, and the graph topology is the structure of the causal graph itself.


This paper reports work at the design and proof-of-concept stage: the theoretical integration, sampling design, and failure-DAG pipeline are specified, while the API-querying and full-scale empirical run are in progress as part of an ongoing undergraduate thesis at Facultad de Ciencias, UNAM. We refrain from reporting preliminary numerical results to avoid presenting underpowered figures; the contribution at this stage is the framework and protocol, offered for discussion ahead of the empirical phase.

\section{Relevance to Causal Discovery and Conclusion}

This work bears directly on the growing practice of using LLMs as sources of expert-like background knowledge in hybrid causal discovery pipelines \cite{long2023can,verma2025causal,butkus2025causal,zhang2023understanding}. If most of an LLM’s apparent causal ability is just associational, then priors extracted from it may work well for tasks that mainly reflect how often concepts co-occur in its training data; but should not be trusted, without separate, external confirmation validation, for orienting edges or resolving Markov-equivalence classes, tasks that are intrinsically interventional or counterfactual. Framing the diagnosis in terms of Pearl's hierarchy yields a criterion portable across models and benchmarks, and the failure-DAG visualization offers a reusable auditing tool for characterizing \emph{where} a model's causal competence breaks down before it is used as a knowledge source. Future work will complete the empirical run across multiple LLM families, report results by level and subtype, and examine whether failure DAGs generalize across models, informing guidance for the use of LLM-elicited priors in causal discovery.

\begin{editorcomment}
Comment from editor: strongly suggest including a figure with a toy example of a “failure DAG” (step iv of the planned pipeline). It is hard to visualize how one of these graphs would look like.
\end{editorcomment}


\begin{editorcomment}
Comment from editor: As this is a proof of concept and the project is in an early stage, please discuss more thoroughly the limitations and challenges that you are already confronting, or you expect to confront, and if possible how they could be addressed
\end{editorcomment}


\section{Background}


Pearl plantea como partida tres habilidades distintas de la cognición: ver, hacer e imaginar
o hipotetizar/anticipar(Rosen). \cite{pearl2018bookOfWhy}
\begin{itemize}
    \item \textbf{Ver - Imitar} Nos ayuda a percibir regularidades las observaciones. Estas
Rosen las pondría como la identificacion de perceptos. 
    \item \textbf{Hacer - Planear} Predecir los efectos provocados por cierta perturbacion en el ambiente
y tomar una decision para obtener el resultado deseado. Esto incluye la habilidad de
tomar herramientas intencionalmente para lograr el objetivo. 
    \item \textbf{Imaginar - Hipotetizar} Cuando tenemos una teoría de porque funciona una herramienta o cuando no, 
podemos hipotetizar y mejorar estas herramientas para futuros escenarios. 
\end{itemize}

Si una de las intenciones de Turing era intentar definir un agente coginitvo en cuanto a
las preguntas que puede responder. La intencion de Pearl aqui plantea, en cada uno de los 
niveles, distintas dificultades. Para Pearl un sistema que realmente presente cognicion 
deberia ser posible de reproducir estos tres tipos de inferencias causales. 

Los tres puntos anteriores representan entonces, la intuicion por detras de la 
formulacion matematica, de dichas relaciones. 

\begin{itemize}
    \item L1. \textbf{Asociativo}. ¿Que pasa si veo ? ¿Como cambiaria mi creencia en Y viendo X? En ejemplos 
    de la vida real sería que me dice un sintoma acerca de una enfermedad ? que me dice el resultado de una encuesta para predecir 
    el resultado final de las elecciones ? 
    \item L2. \textbf{Intervencion}. Que pasaria si hago \dots ? Como seria Y si hago X ? Como puedo hacer que pasae Y ?
    \item L3. \textbf{Contrafactual}. ¿Si tomo una aspirina, se me quitaria mi migraña ?¿Que pasa si prohibimos los cigarros?
\end{itemize}

¿Que es lo que puede hacer un razonador causal?¿Que puede computar un organismo con un modelo causal, 
a falta de el, que le permite predecir? 

Aqui Rosen diferia, en cuanto a que el proceso de inferencia de un organismo, 
no de depende solo de su estructura causal sino tambien viene desde una idea que 
engloba el completo funcionamiento del organismo, definiendo la coginicion y agencia del 
organismo no solo en la estructura causal del fenomeno que se esta tratando de predecir. 

Uno de los objetivos principales de la version fuerte de IA, es reproducir la inteligencia humana. En lugar
de tener una inteligencia genuina, nos encontramos con habilidades impresionantes. La diferencia recae en la falta 
de un modelo de la realidad. 

Muchos de los sistamas de IA se basan especialmente en la asociacion. Toman una colección gigante de observaciones para 
encajar una funcion de la misma manera que un estadista podria querer encajar una linea para pasar por una 
coleccion de puntos. Esto pone una barrera, para que los sistemas puedan ser capaces de reaccionar a circunstancias fuera del 
dataset. Como en los coches de manejo automatico, ¿que pasa si nos encontramos por accidente una vaca en la carretera, un bache
o un borracho ?¿Como podemos hacer un deduccion generalizada? Esta falta de flexibilidad y adaptacion es ineblitable 
en la primerca jerarquia causal. 

\begin{editorcomment}
Comment from editor: The authors cited other papers exploring the relationship between LLMs and causal outputs, it would be relevant to include more detail of how these papers approached the questions pursued by the authors in this project.
\end{editorcomment}



\section{CLADDER}

Causal Inference Engine. 

Data sets are in the form of $D := \{q_i, a_i, e_i \}$, where $q_i$ is a query, $a_i$ is the answer, and $e_i$ is the explanation. The experiment consist in translating the query $q_i$ into a prompt $p_i$ and then sending it to the LLM, $f : q_i \rightarrow p_i$. All the statements contain a Formal part, and a natural language part. 

The LLM will return an answer $a'_i$ and an explanation $e'_i$. The answer $a'_i$ is compared with the ground truth answer $a_i$ to determine if the LLM answered correctly or not. The explanation $e'_i$ is compared with the ground truth explanation $e_i$ to determine if the LLM explained correctly or not. The results are stored in a data frame for further analysis.

All the queries contain statements in the three rungs, and are constructued with binary variables. Each query contain 
three or four nodes to mitigate the lack of mathmatical and large processing mistakes LLMs can make. 

The Causal Cot promt divides the queries in this steps, inspired by 
making a CI engine plugin to aument the LLMs capabilities. The workflow consist in the following steps:

\begin{itemize}
    \item causal graph extration
    \item correct query type interpretation 
    \item symbolic formalization of the query 
    \item semantic parsing to compile the available data
    \item stimand derivation 
    \item arithmetic calculation 
\end{itemize}

As the statements where algorithmically generated, gramatical errors where analyzed using the LanguageTool API, 1.26 grammatical errors per 100 words (98.74\%  accuracy).

Other research effords have focus in checking causality as knowladge, causality as language comprehension, 
and causality as reasioning. 

Other examples of commonsense causality includes 
\begin{itemize}
    \item cause and effect of an agent action 
    \item motivation and emotional reaction in social context 
    \item correspondance of a set of steps with a high level goal
    \item development of a story given a different beginning 
    \item knowladge base for causality 
    \item infer causation from correlation 
\end{itemize}

Remark the importance of causality in policy, desicion making, epidemiology, economics, fairness, and explicability. 

\section{Experimental Setup}

Using Gemma4, Qwen3.6-35B-A3B, Nemotron 3 Super 120B-A12B, Mistral Medium 3.5 128B, DeepSeek-V4-Flash. Open source LLMs, and GPT-4. To compare the performance of the LLMs, we will use the CLADDER benchmark, which contains 10,000 queries in the three rungs of Pearl's ladder. The queries are divided into three levels: L1, L2, and L3. Each level contains a set of queries that test the LLM's ability to reason about causality. 

To mitigate data contamination, we will translate the queries into Spanish, and then back into English using a different LLM. This will ensure that the queries are not present in the training data of the LLMs being tested.  



%
% ---- Bibliography ----
%
\bibliographystyle{splncs04}
\bibliography{references}
\end{document}
